\section{Introduction}\label{sec:introduction}
Let \(\omega_0\) be a smooth K\"ahler form on a compact complex manifold \(M\).
We use the scalar-curvature convention
\(S(\omega)=\operatorname{tr}_\omega\operatorname{Ric}(\omega)\).
The Calabi flow is
\begin{equation}\label{eq:calabi-flow}
 \omega(t)=\omega_0+i\partial\bar\partial\phi(t)>0,
 \qquad \phi(0)=0,\qquad
 \partial_t\phi=S(\omega(t))-\bar S,
\end{equation}
where \(\bar S\) is the average scalar curvature, determined by the K\"ahler class.
Calabi introduced this flow to seek canonical metrics within a K\"ahler
class~\cite{Calabi1982}. Chen's smooth long-time existence conjecture
asserts that the flow exists for all positive times from arbitrary smooth
K\"ahler initial data; see \cite[Conjecture~1.1]{HeZeng2021}.
We give a counterexample on projective space. Write
\(\omega_{\mathrm{FS}}\) for the Fubini--Study form whose potential
on the standard affine chart is \(\log(1+|Z|^2)\).

\begin{theorem}\label{thm:main}
There is a smooth \(U(10)\)-invariant K\"ahler form \(\omega_0\) on
\(\CP^{10}\), in the class \([\omega_{\mathrm{FS}}]\), whose maximal smooth
Calabi flow has existence interval \([0,T_*)\) with \(0<T_*<\infty\).
There are a point \(p\in\CP^{10}\) and a constant \(a>0\) such that
\[
 S(\omega(t))(p)\sim a(T_*-t)^{-1/2}\qquad\text{as }t\uparrow T_*.
\]
\end{theorem}

The Fubini--Study metric has constant scalar curvature. Thus the example
also answers negatively the smooth global-existence question for classes
containing a constant-scalar-curvature K\"ahler metric, emphasized by
Chen and Cheng~\cite[p.~912, following Conjecture~1.3]{ChenCheng2021}.
In particular, no assertion of smooth global existence and subsequent
smooth convergence modulo holomorphic automorphisms can hold for arbitrary
smooth initial metrics in such a class: the obstruction already occurs
at a finite time.

The same singularity persists on products with a stationary factor.
\begin{corollary}[Products with stationary factors]\label{cor:product-singularities}
Let \(\omega(t)\), \(T_*\), \(p\), and \(a\) be as in
Theorem~\ref{thm:main}, and let \((Y,\eta)\) be a compact connected
K\"ahler manifold with \(\eta\) of constant scalar curvature.  With
\(\pi_X\) and \(\pi_Y\) the projections from \(\CP^{10}\times Y\),
the maximal smooth Calabi flow from
\(\Omega_0=\pi_X^*\omega_0+\pi_Y^*\eta\) is
\[
 \Omega(t)=\pi_X^*\omega(t)+\pi_Y^*\eta,
 \qquad 0\leq t<T_*.
\]
For every \(y\in Y\),
\(S(\Omega(t))(p,y)\sim a(T_*-t)^{-1/2}\).
In particular, finite-time smooth Calabi-flow singularities occur in
every complex dimension \(n\geq10\), on
\(\CP^{10}\times\CP^{n-10}\), in a class containing a
K\"ahler--Einstein metric.
\end{corollary}
The corollary is proved after Theorem~\ref{thm:main} in
Section~\ref{sec:compact}.

\subsection{Background and significance}
On compact Riemann surfaces, Chru\'sciel proved global forward existence
and convergence to a constant-curvature metric~\cite[Theorem~2.1]{Chrusciel1991}.
Chen gave a different proof using decreasing functionals and concentration
compactness~\cite{Chen2001}, and Struwe developed a unified treatment of
Calabi and normalized Ricci flows on closed surfaces~\cite{Struwe2002}.
In higher dimensions, Chen and He proved local existence and uniqueness
and established stability near a constant-scalar-curvature
metric~\cite[Theorems~1.2 and~3.2]{ChenHe2008}. He and Zeng developed
short-time theory for rougher initial metrics~\cite{HeZeng2021}.
These higher-dimensional results concern either short times or restricted
initial data.

Curvature estimates give continuation criteria for the smooth flow.
Chen and He proved extension under a uniform Ricci-curvature
bound~\cite[Corollary~5.3]{ChenHe2008}, and Chen and Cheng established
extension under a uniform scalar-curvature bound~\cite[Theorem~1.8]{ChenCheng2021}.
Li, Zhang and Zheng state that a uniform spatial $L^p$ scalar-curvature
bound, measured with respect to the evolving metric, suffices when
$p>\dim_{\C}M$~\cite[Theorem~1.2]{LiZhangZheng2024}; Li and Zhang
state a further criterion involving spacetime scalar and
scalar-Laplacian integrability~\cite[Theorem~1.1]{LiZhang2025}.

On a compact almost homogeneous manifold with two hypersurface ends,
Guan proved exponential smooth convergence of the modified Calabi flow
for initial metrics invariant under a maximal compact subgroup of the
automorphism group~\cite[Theorem~A]{Guan2007}.
The two special orbits of the action used here are instead a point and
the hyperplane at infinity. The point has complex codimension ten, so
the two-hypersurface hypothesis fails.

Streets constructed global weak Calabi flows by minimizing
movements~\cite{Streets2014}. Berman, Darvas and Lu identify these flows
with trajectories in a space of finite-energy K\"ahler potentials and show
that, when the class contains a constant-scalar-curvature metric, every
weak trajectory converges to a smooth constant-scalar-curvature potential
in the $d_1$ distance, the extension of the $L^1$ path-length distance on
smooth K\"ahler potentials~\cite[Section~6 and Theorem~1.11]{BermanDarvasLu2017}.
The weak and classical flows agree for as long as the latter exists
smoothly~\cite[Proposition~6.1]{BermanDarvasLu2017}.
Neither this weak continuation nor the large-time $d_1$ convergence implies
smooth continuation through every finite time.

\subsection{Structure of the construction}

The construction first reduces the flow to a scalar equation, then builds
an exact shrinking solution of its flat limit, and finally corrects that
solution on the compact manifold. The geometric starting point is the
radial ansatz of Calabi~\cite{Calabi1979} and the momentum formulation
of Hwang and Singer~\cite{HwangSinger2002}. Abreu's toric
scalar-curvature formula~\cite{Abreu1998} is expressed in Guillemin's
symplectic-potential framework~\cite{Guillemin1994}.
For a radial potential \(\Phi(\ell)\), where
\(\ell=\log|Z|^2\), set \(x=\Phi'(\ell)\) and
\(v=\Phi''(\ell)\). We use \(x\in(0,1)\) as the momentum
coordinate and write
\[
 v=x(1-x)+x^2(1-x)^2H.
\]
Smoothness and positivity at both special orbits are encoded by
\(H\in C^\infty[0,1]\) and \(1+x(1-x)H>0\).
Section~\ref{sec:geometry} derives the scalar fourth-order equation
for \(H\), including the affine choice in the Legendre potential
needed to recover the Calabi flow with its prescribed normalization.
It also realizes the fixed spatial operator on an auxiliary sphere,
where the endpoint conditions become ordinary smooth invariance.

The singularity is detected by one value of \(H\). At the fixed
affine origin \(p\), the scalar-curvature identity is
\[
 S(t,p)=110\bigl(1-H(t,0)\bigr).
\]
We construct a smooth positive profile \(G\), bounded above and away from
zero, with \(G(0)=1\) and \(G'(0)<0\), and set
\[
 F(t,x)=\rho^{-1}h(x/\rho),\qquad
 h(y)=\frac{G(y)-1}{y},\qquad
 \rho=\sqrt{-t},\quad t<0.
\]
Then \(F(t,0)=\rho^{-1}G'(0)\). Our task is to find an exact
compact solution \(H=F+V\) for which positivity persists and
\[
 V(t,0)=O(\rho^{-1+\delta})\qquad\text{for some }\delta>0.
\]
The correction is therefore smaller than the leading term at the
origin, preserving the desired scalar-curvature divergence. A smooth
interior time slice of this solution will be the initial metric;
no initial metric needs to be fixed before the construction.

Section~\ref{sec:profile} constructs \(G\) so that \(F\) solves the
scale-invariant flat equation exactly. The profile is obtained by
matching a two-parameter analytic solution at the origin to a half-line
solution approaching a positive constant.
A finite exact-arithmetic calculation, with analytic enclosure and
nonlinear error bounds, certifies the matching inequalities. This uses
the framework of validated Taylor integration surveyed
in~\cite{NedialkovJacksonCorliss1999}.
The finite argument and its complete executable certificate appear in
the appendices. The resulting profile also has a differentiable
expansion at infinity; this expansion controls the compact error away
from the shrinking core.

The correction begins with a right inverse for the flat linearized
equation. Section~\ref{sec:parabolic} introduces the weighted norms
and local estimates, and Section~\ref{sec:flat} constructs this inverse
on all negative times. In similarity variables its time-one evolution
is a compact perturbation of a damped constant-coefficient evolution.
A spectral split determines finitely many components from future
forcing. This choice is available because the initial metric is still
free; it requires no stability or nondegeneracy theorem for the
profile. The compact-perturbation and spectral-projection tools are
classical~\cite{Kato1995}. After selecting a spectral weight, an outer
annulus estimate upgrades the temporal supremum bound to the full
spatially weighted parabolic H\"older estimate. The decisive input
there is a scale-invariant \(L^2\) bound that is integrable up to
the singular time.

Section~\ref{sec:compact} combines this ancient inverse near the core
with a short-time forward solution outside it. The outer forcing
cancels the inner cutoff commutator exactly; the remaining operator
errors are small. For the nonlinear correction one further needs the
compact residual to become small after applying the inverse. This
holds in a norm controlling all four spatial derivatives and their
temporal H\"older seminorms, with the time derivative omitted.
The outer estimate uses a strictly higher H\"older exponent for the
source, and the fixed-point equation then supplies the missing time
derivative. Positivity and the pointwise correction bound above give
the geometric hypotheses needed to finish Theorem~\ref{thm:main}.
